
% Exposé VII: back to the volume's Sym and Proj
\renewcommand{\Sym}{\operatorname{Sym}}
\renewcommand{\Projz}{\operatorname{Proj}_0}
\section*{Exposé VII\\[3pt] IMMERSIONS RÉGULIÈRES ET CALCUL DU \texorpdfstring{@@M0@@}{K0} D'UN SCHÉMA ÉCLATÉ}
\begin{center}
\uline{par Pierre BERTHELOT}
\end{center}

\section*{1. Généralités sur les immersions régulières}

La notion classique d'immersion régulière (EGA IV 16.9.2) s'avérant inadéquate dans le cas non noethérien, nous allons définir une notion plus faible, mais donnant lieu aux propriétés souhaitées, et qui équivaut à la définition usuelle dans le cas noethérien.

\textbf{Définition 1.1.} \uline{Soient} @@M1@@ \uline{un anneau commutatif}, @@M2@@ \uline{un} @@M3@@-\uline{module projectif de type fini}, @@M4@@ \uline{un homomorphisme} @@M5@@-\uline{linéaire}. \uline{On dit que} @@M6@@ \uline{est régulier si le complexe de Koszul} @@M7@@ \uline{défini par} @@M8@@ \uline{est acyclique en dimension} @@M9@@, \uline{donc est une résolution de} @@M10@@.

Rappelons que le complexe de Koszul associé à @@M11@@ est défini comme suit: le @@M12@@-module gradué @@M13@@ est l'algèbre extérieure @@M14@@, et l'opérateur bord la multiplication intérieure par @@M15@@ considéré comme élément du dual @@M16@@ : c'est une antidérivation de degré @@M17@@ de @@M18@@, égale à @@M19@@ sur @@M20@@.

\textbf{Lemme 1.2.} \uline{Soit} @@M21@@ \uline{un homomorphisme régulier}. \uline{Soient} @@M22@@ \uline{une} @@M23@@-\uline{algèbre}, @@M24@@ \uline{l'homomorphisme déduit de} @@M25@@ \uline{par changement de base}; \uline{on suppose} @@M26@@ \uline{régulier}. \uline{On pose} @@M27@@, @@M28@@. \uline{Alors on a les relations suivantes pour tout entier} @@M29@@:
@@M30@@

a) On suppose d'abord que les @@M31@@ sont projectifs sur @@M32@@; comme @@M33@@ est de type fini, ils sont localement libres de type fini sur @@M34@@.
